(1) 定义$x=x_0(a+a^\dagger)$，$x_0=\sqrt{\hbar/(2m\omega)}$，并用$[a,a^\dagger]=1$正规排序。$x^4$作用于$|n\rangle$后只含$|n-4\rangle,|n-2\rangle,|n\rangle,|n+2\rangle,|n+4\rangle$分量，因而纯偶宇称耦合使$\Delta n=\pm1,\pm3$矩阵元为零。这并不是偶极辐射跃迁选择定则，而是四次势本身的算符奇偶性。
(2) 由$a|n\rangle=\sqrt n|n-1\rangle$和$a^\dagger|n\rangle=\sqrt{n+1}|n+1\rangle$逐次作用，取$|n\rangle$的系数得$\langle n|(a+a^\dagger)^4|n\rangle=6n^2+6n+3=3(2n^2+2n+1)$。因此$E_n=\hbar\omega(n+\tfrac12)+\frac{3\lambda\hbar^2}{4m^2\omega^2}(2n^2+2n+1)+O(\lambda^2)$，基态修正为$3\lambda\hbar^2/(4m^2\omega^2)>0$。
(3) 两相邻能级之差$\Delta E_{n+1,n}=E_{n+1}-E_n=\hbar\omega+\frac{3\lambda\hbar^2}{m^2\omega^2}(n+1)+O(\lambda^2)$，故跃迁角频率$\omega_{n+1,n}=\Delta E/\hbar=\omega+\frac{3\lambda\hbar}{m^2\omega^2}(n+1)+O(\lambda^2)$。四次正势导致能级随$n$愈发疏远，等时间隔的理想谐振子性质被破坏。
(4) 高$n$时一阶修正量级为$\lambda\hbar^2 n^2/(m^2\omega^2)$，未扰动能量间距仍为$\hbar\omega$。确保相邻态混合受控的充分小参数约为$\lambda\hbar n^2/(m^2\omega^3)\ll1$（更精细地需检查非对角矩阵元与$2\hbar\omega$、$4\hbar\omega$间隙之比）。固定$n$令$\lambda\to0$时全部修正消失，严格回到等间距能谱。
